% =====================================================================
% §2.7 — La molteplicità delle voci (esempio unico: PGE_voices)
%
% Revisione: un solo esempio (PGE_voices.yml) con scatter DENTRO, dopo
%   l'esempio completo, come estensione. La MAP si descrive per fatti:
%   - num_voices envelope (50 -> 7 -> 70): molte voci, asciugate al centro,
%     ricomposte in chiusura; la densità segue circa lo stesso andamento;
%   - distribution: loop nel primo 40%, poi asincrono pieno (1) verso il centro;
%   - scatter: apre solo nella seconda metà -> decorrelazione temporale fra
%     le voci, leggibile a bassa densità;
%   - parametri spaziali del blocco voices (pan strategy) correlati a
%     pan/pan_range della singola voce; non sugli assi della MAP.
% =====================================================================

% ----------------------------------------------------------------
\subsection{La molteplicità delle voci}\label{sec:voci}

Tutto il percorso fin qui ha riguardato una voce sola. Le voci non sono
sorgenti indipendenti: leggono tutte lo stesso materiale e ripetono lo stesso
gesto, e la coerenza è garantita per costruzione --- la voce~0 è il
riferimento, e ogni asse di differenziazione misura uno scarto rispetto a
essa, nullo per definizione sulla voce~0. Ciò che resta aperto è
l'\emph{ampiezza} di quello scarto, e con essa il risultato percettivo. Sotto
la soglia dei pochi millisecondi gli scarti restano attributo morfologico e
la massa si ispessisce; oltre quella soglia diventano eventi
separati~\cite{Vaggione2002}, e lo stesso meccanismo produce polifonia. È il
regime in cui lavora Truax quando assegna a ciascuna voce un armonico
indipendente sopra una fondamentale comune, e insieme sovrappone fino a
diciotto voci non sincronizzate sullo stesso campione~\cite{Truax1994}: la
scelta di dove collocarsi lungo quel continuo è una decisione compositiva, e
si dichiara nella misura dei passi.

Gli assi sono quattro, ortogonali, ciascuno governato da una strategia
intercambiabile dichiarata per nome. La trasposizione (\texttt{pitch}) è un
fattore moltiplicativo sul rapporto di lettura: a gradini regolari, spalmata
su un intervallo dato, per gradi di un accordo nominale, o estratta a sorte
una volta per voce. La posizione di lettura (\texttt{pointer}) sposta ciascuna
voce nel buffer, a passi equidistanti o per scarti casuali attorno alla
traiettoria comune. L'attacco (\texttt{onset}) ritarda l'ingresso della voce,
con passo lineare o esponenziale: gli scarti sono sempre positivi, le voci
seguono e non precedono mai. Il pan distribuisce le voci nel campo stereo.
Nessuno di questi passi è obbligato a restare fisso: dove l'esempio scrive
\texttt{step} può comparire un inviluppo, e la differenziazione fra le voci
si apre o si chiude nel corso del brano.

Due inviluppi governano poi l'insieme. Il primo è il numero di voci:
\texttt{num\_voices} scende e risale nel corso del brano, e con esso il numero
di grani in sovrapposizione --- meno voci attive, meno grani. Nell'esempio
(Fig.~\ref{fig:voci-MAP}) la posizione di lettura è distribuita a passi
regolari, e la trasposizione procede a gradini di un'ottava per voce: uno
scarto largo, lontano dal regime della decorrelazione, che distende le voci
su un registro vastissimo --- la scala cromatica della figura misura
l'escursione risultante.

Il secondo governa l'accoppiamento temporale, ed è il livello che si aggiunge
alla griglia della singola voce (\S\ref{sec:griglia}): il ritardo fra i grani
di un flusso e il grado di sincronicità fra flussi diversi sono due controlli
distinti, e la stessa coppia struttura per intero il trattamento canonico
della granulazione~\cite{Dutilleux2016}. Qui, su un fondo di griglia
pienamente asincrona (\texttt{distribution} a~1), agisce \texttt{scatter}: a~0
tutte le voci condividono la sequenza di intervalli della voce~0, a~1 ciascuna
calcola la propria, i valori intermedi interpolano. Si distingue dallo scarto
di attacco, che è uno sfasamento fisso: \texttt{scatter} non ritarda le voci,
decide quanto la loro scansione temporale sia la stessa. Anche questo calcolo
è stocastico: quali voci si sganciano e in quali istanti cambia a ogni
esecuzione, ma l'inviluppo di \texttt{scatter} resta invariante.

Restano i parametri spaziali, dichiarati nel blocco ma non riportati sugli
assi della \textsc{map}. La spazializzazione agisce a due livelli annidati:
la strategia \texttt{pan} delle voci assegna a ciascuna un proprio angolo,
distribuendole nel campo --- e qui il passo è un inviluppo, sicché il
ventaglio resta chiuso per metà brano e si apre nell'ultimo terzo --- mentre
il \texttt{pan\_range} dello stream aggiunge a ogni grano della singola voce
una deviazione spaziale attorno a quell'angolo.

Le quattro famiglie di strategie sono un punto di estensione dichiarato:
aggiungere un modo diverso di distribuire le voci lungo uno degli assi
significa scrivere una classe che ne calcoli lo scarto e registrarla sotto un
nome, che diventa immediatamente scrivibile nel listato. Ciò che invece resta
un intervento di struttura è l'aggiunta di un asse nuovo, che non si esaurisce
in una classe.

\begin{figure}[H]
    \begin{framedfloat}
    \lstinputlisting[
      linerange={47-61},
      language=yaml,
      basicstyle=\ttfamily\fontsize{7}{8.4}\selectfont
    ]{examples/voices/PGE_voices.yml}
    \vspace{2pt}
    \includegraphics[width=\linewidth]{examples/voices/PGE_voices_map}
    \caption{In alto, le chiavi salienti di \texttt{PGE\_voices.yml}:
    \texttt{num\_voices} e \texttt{scatter} come inviluppi, la
    \texttt{distribution} con loop locale e le strategie del blocco
    \texttt{voices}. Sotto, la \textsc{map} (assi come in
    Fig.~\ref{fig:c-e}). Ogni voce traccia una banda orizzontale sull'asse
    della posizione di lettura; il colore codifica la trasposizione, distinta per
    voce. Le bande sono fitte e dense di grani in apertura, si diradano verso il
    centro --- dove il numero di voci è minimo e i singoli onset diventano
    leggibili --- e si moltiplicano di nuovo in chiusura. Nella seconda metà le
    bande, inizialmente allineate sulla stessa griglia temporale, si sfasano
    progressivamente fra loro fino a confondersi in una nuvola diffusa.}
    \label{fig:voci-MAP}
    \end{framedfloat}
\end{figure}
